
This study examines whether code generation models aligned with different training objectives produce different error type distributions among failed generations. Two models were compared: CodeRL (770M parameters, trained with execution-based reinforcement learning) and CodeLlama-7B-Instruct (7B parameters, instruction-tuned). On a combined dataset of HumanEval+ and MBPP+ benchmarks (542 problems, 766 total failures), error type distributions differed between the two models (chi-square = 35.27, p = 2.19 $\times 10^{-8}$, Cramér's V = 0.21). Among RL model failures, 2.12\% were assertion errors compared to 0\% for the instruction-tuned model (Fisher's exact test, p = 0.0027). Using execution tracing, RL-generated code executed a mean of 29.4\% of lines before failure compared to 0.09\% for the instruction-tuned model (Welch's t-test, p = 1.08 $\times 10^{-34}$, Cohen's d = 1.69). The effect size increased from Cramér's V = 0.21 at the coarse 3-category taxonomy to V = 0.82 at the fine-grained 19-category LlmFix taxonomy. These results are confounded by differences in model architecture (encoder-decoder vs. decoder-only), model size (770M vs. 7B parameters), and pre-training data. The findings do not establish that alignment method is the causal factor for the observed differences.
